\ansclass{10}{Observability}
\ans{1} (a) \final{(1) metrics, (2) tracing, (3) logs, (4) alerting, (5) health check} (readiness). (b) Rate, 5xx ratio and p99 of \texttt{POST /checkout} are \final{RED}; database CPU and the pool's pending acquisitions are \final{USE}. \final{Pending acquisitions is the saturation signal}: work waiting for a connection, and it rose before the p99 did. (c) CPU and the global error rate are \emph{causes} averaged over every route and pod. Fans waiting on a pool use no CPU and do not fail (they wait), so the symptom is only visible in the \final{latency SLI of the checkout journey} and in saturation.
\careful{watching resource dashboards (CPU, pods up) and calling them ``the health of the system''. Users feel latency and errors per journey.}
\ans{2} Series multiply: \final{17; 68; 1\,700; 10\,200; 204\,000; $2.04 \times 10^{11}$}. Memory at 4 KiB each: 68 KiB; 272 KiB; 6.6 MiB; 39.8 MiB; \final{$\approx$ 797 MiB}; \final{$\approx$ 760 TiB}. (b) The five bounded labels give about 200 thousand series, which a metrics store handles; \texttt{customer\_id} multiplies them by a million. Per-customer detail belongs \final{in span attributes and log fields} (searchable by \texttt{trace\_id} or order id), never in a metric label.
\careful{judging a label by its own size (``it is just one more label'') instead of multiplying it by all the others.}
\ans{3} (a) $(180 + 180 + 180 + 1\,260)/4 =$ \final{450 ms}. (b) 1\% of 100\,000 is \final{1\,000 requests}, almost all from pod \texttt{d}: its slowest 10\% (1\,000 of its 10\,000) are slower than anything on the healthy pods. So the fleet p99 $\approx$ the p90 of pod \texttt{d} $\approx$ \final{760 ms} (the class computes 756 ms). The dashboard understates it by 300 ms. (c) A 0.3 s request increments \final{\texttt{le=0.5}, \texttt{le=1}, \texttt{le=2.5} and \texttt{+Inf}}. \final{No}: without a bucket at 2 s the share under 2 s can only be bracketed between the 1 s and 2.5 s buckets; add \texttt{le=2} (the SLO threshold must be a bucket edge). (d) \final{Sum the bucket counters of all pods per \texttt{le}, then compute the 99th percentile once} (\texttt{histogram\_quantile(0.99, sum by (le) (rate(\dots)))}).
\careful{averaging percentiles. The average ignores how many requests each pod served and where the tail is.}
\ans{4} (a) Traces $500 \times 8 \times 500$ B $=$ \final{2 MB/s $\approx$ 173 GB/day}, over \final{5 TB per 30 days}; logs $500 \times 5 \times 300$ B $=$ \final{0.75 MB/s $\approx$ 65 GB/day}. (b) Everything: \final{691 GB/day}, 1\,728\,000 error traces. Head 10\%: 69 GB, \final{only 172\,800 error traces} (90\% lost). Tail: $0.01 + 0.02 + 0.05 \times 0.97 = 7.85\%$ of the traces, \final{54 GB/day with all 1\,728\,000 error traces}. (c) \final{Tail sampling}: it keeps exactly the traces you will open in an incident and scales the rest down. Rates, error ratios and percentiles must come from all requests, otherwise the SLO and the alerts would be computed on a biased sample.
\careful{head sampling at 1\% ``to save money'' with no tail sampling: the rare error trace you need is almost never kept.}
\ans{5} (a) $0.001 \times 672 \times 60 =$ \final{40.3 min}. (b) 1.44\%: \final{burn 14.4, $\approx$ 2\% of the budget per hour, gone in 47 h $\approx$ 1.9 days}. 0.6\%: \final{burn 6, $\approx$ 5\% in 6 h, gone in 112 h $\approx$ 4.7 days}. (c) 100\%: $60 \times 0.0144 \approx$ \final{0.9 min}; 2\%: \final{43 min}. At 0.7\% the 1 h ratio never reaches 1.44\%: \final{only the 6 h rule fires}, after $360 \times 0.6/0.7 \approx$ \final{309 min $\approx$ 5 h}; below 0.6\% only the ticket. (d) $60 \times 1.44/4 \approx$ \final{22 min, at about 10:25}: most of the sale is gone. \final{Give the on-sale its own SLO} (e.g. 99\% of checkouts under 2 s during the on-sale hour) with 5-minute and 1-minute windows.
\careful{one alerting setup for every situation. A budget tuned for 28 days cannot protect an event that lasts 30 minutes.}
\ans{6} (a) The critical path is gateway $\to$ \texttt{order} $\to$ \texttt{payment} $\to$ \final{\texttt{pool.acquire}}: $7.90/9.12 \approx$ \final{87\%} of the time is queueing inside BoxOffice, not the provider (1.02 s). (b) $X = L/W = 20/1.02 \approx$ \final{20 charges/s}. Above that, requests queue for a connection; the wait grows while the arrivals exceed 20/s, and every waiting request holds a thread in \texttt{order} and in the gateway. (c) Symptom: checkout p99 9 s with CPU 45\%. Cause: the provider got slower, each call held its connection longer and the pool of 20 saturated (Little's Law); waiting uses no CPU and shallow health checks still pass. Fix: \final{a bounded wait on the pool (e.g. 500 ms) and a timeout on the provider call} (plus a bulkhead and a fallback, class 9); watch \final{the pool's pending acquisitions} and add a span for the wait.
\careful{blaming the provider for the whole 9 s. Without the \texttt{pool.acquire} span the trace points at the wrong team.}
\ans{7} (a) \final{MTTD 1 min, MTTA 1 min, MTTM 12 min, MTTR 38 min}. (b) $200 \times 720 \text{ s} \times 0.12 =$ \final{17\,280 failed checkouts}; the 28-day budget is $2\,000\,000 \times 0.001 = 2\,000$, so the incident spent \final{$\approx$ 8.6 windows of budget (about 8 months) in 12 minutes}. (c) A good drawing: every service (gateway, \texttt{order}, \texttt{payment}) and the database exports RED/USE to Metrics; the \texttt{trace\_id} enters at the gateway and flows through \texttt{traceparent} (and Kafka headers) to Tracing and into every log line; Health Checks on the load balancer use local readiness only; Alerting reads the checkout SLI from Metrics, \final{pages on the fast burn and on the on-sale SLO, opens a ticket on the slow burn}, with a runbook and a dead man's switch. (d) For example: \final{timeout of 800 ms $+$ bulkhead on the payment adapter} (prevent, payments team); \final{on-sale SLO with 1 min and 5 min burn alerts} (detect, SRE); provider latency in the pre-sale load test (prevent, performance). Each with an owner and a date.
\careful{action items like ``be more careful'' with no owner and no date; they never get done.}
